\section{Experimental Setup}
\label{sec:exp}

\subsection{Data Collection}

We collect prediction market data from Polymarket via two complementary APIs. The Gamma Markets API provides event metadata including titles, descriptions, category labels, creation timestamps, and resolution outcomes. The Central Limit Order Book (CLOB) API supplies real-time and historical price data, order book depth, and trading volume. Together, these sources provide comprehensive coverage of market activity.

Our dataset comprises three tiers of increasing temporal specificity:

\textbf{Historical corpus.} We collect 1{,}443 resolved events across all available categories, spanning January 2024 through March 2026. Each event record includes the final resolution (binary outcome), the market price at various time points before resolution, cumulative trading volume, and metadata sufficient for domain classification. This corpus serves as the primary evaluation set for backtesting and evolution.

\textbf{Batch predictions.} To assess calibration with statistical power, we generate predictions for 1{,}032 active markets in a single batch run, recording both our predicted probabilities and the market prices at the time of prediction. As these markets resolve over the following weeks, we accumulate ground-truth labels for calibration curve construction.

\textbf{Live tracking.} We maintain continuous predictions on 15 high-value markets selected for domain diversity and analytical interest, including commodity price targets (crude oil at various strike prices), cryptocurrency milestones (Bitcoin price thresholds), and entertainment awards (Academy Awards Best Picture). These predictions are refreshed with current context at each evaluation cycle and compared against final outcomes upon resolution.

\subsection{Domain Taxonomy}

We classify all events into eleven domains based on keyword matching augmented with LLM disambiguation. Table~\ref{tab:domains} summarizes the distribution and characteristics of each domain in our historical corpus.

\begin{table}[t]
\centering
\caption{Domain taxonomy with event counts and estimated market efficiency. Efficiency ratings reflect the degree to which market prices incorporate available information.}
\label{tab:domains}
\small
\begin{tabular}{@{}llcc@{}}
\toprule
Domain & Events & Efficiency & Predictability \\ \midrule
Sports & 400+ & Very High & Low \\
Crypto & 100+ & Medium & Medium \\
Commodities & 50+ & Medium & Medium-High \\
Economy & 40+ & Medium-High & Medium \\
Politics & 50+ & Medium & Medium \\
Awards & 20+ & Low & High \\
Geopolitics & 30+ & Medium & Low-Medium \\
Technology & 20+ & Low & Medium \\
Entertainment & 15+ & Low & Medium \\
Golf & 30+ & Very High & Low \\
Other & 50+ & Variable & Variable \\
\bottomrule
\end{tabular}
\end{table}

\subsection{Evaluation Metrics}

\textbf{Brier Score.} Our primary metric, defined as the mean squared error between predicted probability $p_i$ and binary outcome $o_i \in \{0, 1\}$:
\begin{equation}
    \text{BS} = \frac{1}{N} \sum_{i=1}^{N} (p_i - o_i)^2
\end{equation}
Lower scores indicate better calibration. We report Brier Score both overall and decomposed by domain to identify where prediction alpha concentrates.

\textbf{Beat-Market Rate.} The fraction of predictions where our Brier Score for a single event is lower than the market's Brier Score, computed as:
\begin{equation}
    \text{BMR} = \frac{1}{N} \sum_{i=1}^{N} \mathbb{1}\left[(p_i - o_i)^2 < (m_i - o_i)^2\right]
\end{equation}
where $m_i$ is the market price at prediction time. A BMR above 50\% indicates systematic outperformance of market consensus.

\textbf{Calibration Curve.} We bin predictions into deciles and compare the mean predicted probability against the observed frequency of positive outcomes within each bin. Perfect calibration produces a diagonal curve; systematic biases manifest as consistent deviations.

\textbf{Simulated Kelly Return.} Using half-Kelly position sizing, we compute the cumulative return that would result from betting on events where our prediction diverges from the market price by more than 3\%. This metric translates probabilistic accuracy into economic terms.

\subsection{Baselines}

We compare Murmur against four baselines:

\begin{enumerate}
    \item \textbf{Market Consensus}: Always predict the current market price. This represents the ``wisdom of crowds'' baseline and is theoretically hard to beat in efficient markets.
    \item \textbf{Zero-Shot LLM}: Llama~4 Scout (17B parameters) prompted with event title and description only, without context injection or market anchoring. This isolates the raw forecasting capability of the LLM.
    \item \textbf{Context-Enriched LLM}: Full Context Planner pipeline with real-time data injection, but without evolutionary optimization of the strategy. This isolates the contribution of context enrichment.
    \item \textbf{Evolved Strategy}: The best strategy discovered by the Evolution Engine, representing pure code-evolved prediction without LLM reasoning at inference time. This isolates the contribution of evolutionary strategy discovery.
\end{enumerate}

\subsection{Implementation Details}

All LLM inference uses Llama~4 Scout 17B via the Groq API, chosen for its favorable cost-latency profile (\$0.001 per prediction, 3--5 seconds per round). Real-time price data is sourced from CoinGecko (cryptocurrency) and Yahoo Finance (commodities, equities, indices). News retrieval uses the Brave Search API with a maximum of three queries per event. The Evolution Engine runs for 20 rounds per session with single-change constraints and a banned modification list. All experiments are conducted on an NVIDIA DGX Spark (ARM64, 128GB unified memory).

Automated evaluation runs on a cron schedule (8:00 AM and 8:00 PM Eastern Time daily), checking for newly resolved markets and updating performance tracking. The full pipeline---from event retrieval through context enrichment, prediction, and logging---processes approximately 1{,}000 events in under 30 seconds.
